\titledquestion{Brownian motion as a Gaussian process}
% @ASSUMPTION_GLOBAL
Let $W_t$ be a Brownian motion.
\begin{parts}
    % @QUESTION
    \part For $t>0$ and $N\in\N^*$, what can one say about the random variable
       $I_N=\frac{t}{N}\sum_{k=1}^NW_{\frac{kt}{N}}$? 
    \begin{xsolution}
    % @LIST_START
    % @ATOM
    % @PRECOND
    $I_N = \frac{t}{N}\sum_{k=1}^N W_{\frac{kt}{N}}$ is a finite linear combination of the Gaussian variables $W_{\frac{kt}{N}}$, and 
    % @PRECOND
    $\mathbb{E}[W_s] = 0$ for all $s$.
    % @PRECOND
    Its mean is $\frac{t}{N}\sum_{k=1}^N \mathbb{E}[W_{\frac{kt}{N}}] = 0$.
    % @ARGUMENT
    A linear combination of Gaussian random variables is Gaussian. 
    % @OUTCOME
    $I_N$ is a centered Gaussian random variable.
    % @LIST_END
    \end{xsolution}
    \part What is the limit if $I_N$ when $N$ goes to infinity?
    \begin{xsolution}
    % @SET_START
    % @LIST_START
    % @ATOM
    % @PRECOND
    For Brownian motion, $\E[W_aW_b]=\min(a,b)$ for all $a,b\ge0$.
    % @PRECOND
    Using bilinearity and the covariance formula,
    \[
    \Var(I_N)
    =\frac{t^2}{N^2}\sum_{k=1}^N\sum_{j=1}^N
    \min\!\left(\frac{kt}{N},\frac{jt}{N}\right).
    \]
    % @ARGUMENT
    This is a double Riemann sum, 
    % @OUTCOME
    hence as $N\to\infty$,
    \[
    \Var(I_N)\xrightarrow[N\to\infty]{} \int_0^t\!\!\int_0^t \min(u,v)\,du\,dv=\frac{t^3}{3}.
    \]
    % @ATOM
    Since 
    % @PRECOND
    $I_N$ is centered Gaussian and 
    % @PRECOND
    its variance converges, 
    % @ARGUMENT
    the distribution is determined by the variance alone.
    % @OUTCOME
    Then in law
    \[
    I_N \xrightarrow[N\to\infty]{}\mathcal{N}\!\left(0,\frac{t^3}{3}\right).
    \]
    % @LIST_END
    % 
    % @LIST_START
    % @ATOM
    % @PRECOND
    $I_N=\sum_{k=1}^N W_{k\Delta t}\,\Delta t$ with $\Delta t=t/N$ is a Riemann sum for $\int_0^t W_s\,ds$. 
    % @PRECOND
    Brownian paths $s\mapsto W_s$ are continuous, hence integrable on $[0,t]$.
    \\
    % @ARGUMENT
    A Riemann sum of a continuous $L^2$ function converges in $L^2$. 
    % @PRECOND <-- I decided to allow precond to not be ordered
    Since each $I_N$ is Gaussian and $L^2$ limits of Gaussian variables are Gaussian, 
    % @OUTCOME
    the limit $I$ is Gaussian with mean $0$ and variance $t^3/3$ in $L^2$.
    % @LIST_END
    % @SET_END
    \end{xsolution}
    % @QUESTION
    \part Deduce that 
        $$X_t = \int_0^t W_s ds$$
        is a centered Gaussian random variable and compute its characteristic function. 
    \begin{xsolution}
    % @LIST_START
    % @ATOM
    % @PRECOND
    From previous question, $I_N \to X_t:=\int_0^t W_s\,ds$ in $L^2$ and 
    % @PRECOND
    each
    \[
    I_N=\frac{t}{N}\sum_{k=1}^N W_{\frac{kt}{N}}
    \]
    is a centered Gaussian random variable. Moreover, 
    % @PRECOND
    $\E[W_s]=0$ for all $s\ge0$, and 
    % @PRECOND
    $\Var(I_N)\to \frac{t^3}{3}$.
    \\
    % end
    % argument
    Since $I_N\to X_t$ in $L^2$, we have $I_N\Rightarrow X_t$ and $\E[I_N^2]\to \E[X_t^2]$.  
    Because $I_N$ is centered Gaussian with $\Var(I_N)\to \frac{t^3}{3}$, we also have
    \[
    I_N \Rightarrow \mathcal N\!\left(0,\frac{t^3}{3}\right).
    \]
    By uniqueness of the limit in distribution, it follows that $X_t$ is Gaussian with the same law. Finally,
    \[
    \E[X_t]=\int_0^t \E[W_s]\,ds=0.
    \]
    % end
    % outcome
    \[
    X_t \sim \mathcal N\!\left(0,\frac{t^3}{3}\right).
    \]
    % end
    
    % step(2)
    % precond
    $X_t \sim \mathcal{N}(0,\sigma^2)$ with $\sigma^2=\frac{t^3}{3}$.
    % end
    % argument
    Use the characteristic function of a centered Gaussian:
    % end
    % outcome
    \[
    \varphi_{X_t}(u)=\E[e^{iuX_t}]
    =\exp\!\left(-\frac12\sigma^2u^2\right)
    =\exp\!\left(-\frac12\cdot\frac{t^3}{3}u^2\right)
    =\exp\!\left(-\frac{t^3}{6}u^2\right)
    \]
    % @LIST_END
    \end{xsolution}
% question
    \part Compute the first 2 moments of the random variable
        $$
        Y:= \int_0^t W_s^2 ds. 
        $$

\begin{xsolution}

% step(1)
% precond
$Y=\int_0^t W_s^2\,ds$ and for Brownian motion $\E[W_s^2]=\Var(W_s)=s$ for all $s\in[0,t]$.
% end
% argument
Since $W_s^2\ge 0$, by Fubini's theorem:
% end
% outcome
\[
\E[Y]=\int_0^t \E[W_s^2]\,ds=\int_0^t s\,ds=\frac{t^2}{2}.
\]
% end

% step(2)
% precond
$Y^2=\left(\int_0^t W_s^2\,ds\right)^2=\int_0^t\!\!\int_0^t W_s^2W_u^2\,ds\,du$ and $W_s^2W_u^2\ge0$.
% end
% argument
By Fubini's theorem (nonnegative integrand), swap expectation and double integral:
% end
% outcome
\[
\E[Y^2]=\int_0^t\!\!\int_0^t \E[W_s^2W_u^2]\,ds\,du.
\]
% end

% step(3)
% precond
For $s,u\in[0,t]$, $(W_s,W_u)$ is centered Gaussian with $\E[W_s^2]=s$, $\E[W_u^2]=u$, $\E[W_sW_u]=\min(s,u)$.
% end
% argument
By the Isserlis--Wick formula for centered Gaussians:
% end
% outcome
\[
\E[W_s^2W_u^2]=\E[W_s^2]\E[W_u^2]+2(\E[W_sW_u])^2=su+2\min(s,u)^2.
\]
% end

% step(4)
% precond
$\E[Y^2]=\int_0^t\!\!\int_0^t \E[W_s^2W_u^2]\,ds\,du$ (Step 2) and $\E[W_s^2W_u^2]=su+2\min(s,u)^2$ (Step 3).
% end
% argument
Substitute and compute the two integrals by symmetry of $[0,t]^2$:
\[
\int_0^t\!\!\int_0^t su\,ds\,du=\left(\frac{t^2}{2}\right)^2=\frac{t^4}{4},
\qquad
\int_0^t\!\!\int_0^t \min(s,u)^2\,ds\,du
=2\int_0^t\!\int_0^u s^2\,ds\,du
=2\int_0^t \frac{u^3}{3}\,du=\frac{t^4}{6}.
\]
% end
% outcome
\[
\E[Y^2]=\frac{t^4}{4}+2\cdot\frac{t^4}{6}=\frac{7t^4}{12}.
\]
% end
    \end{xsolution}
    \end{parts}